% ---------------------------------------------------------------------
\section{Comment ça fonctionne, concrètement}
% ---------------------------------------------------------------------

\subsection{La chaîne, de votre fichier au lecteur}

\begin{enumerate}
  \item \textbf{Vous préparez deux fichiers} : l’intérieur (toutes les
        pages, en PDF) et la couverture (première, dos, quatrième en un
        seul PDF « à plat »).
  \item \textbf{Vous les déposez sur une plateforme} (KDP, BoD,
        IngramSpark, Bookelis…), avec les métadonnées : titre, auteur,
        ISBN, prix, résumé, catégories, mots-clés.
  \item \textbf{La plateforme contrôle les fichiers} et vous propose une
        épreuve numérique, puis papier (fortement recommandée).
  \item \textbf{Le titre est mis en vente} : sur la boutique de la
        plateforme, chez les libraires en ligne et, selon les offres,
        dans les bases professionnelles consultées par les librairies
        (Dilicom, Electre).
  \item \textbf{À chaque commande}, l’exemplaire est imprimé, relié,
        emballé et expédié. Vous touchez votre redevance ou votre marge,
        généralement chaque mois, avec un décalage de 30 à 90 jours.
\end{enumerate}

\subsection{Les machines}

L’IAD repose sur des presses numériques (toner ou jet d’encre haute
vitesse) couplées à des lignes de finition automatisées : assemblage,
reliure, massicotage trois faces. Un livre peut passer du fichier à
l’exemplaire relié en quelques minutes. Les principaux sites qui
impriment pour le marché français se trouvent en France, en Allemagne,
au Royaume-Uni, en Pologne, en Italie ou en Espagne, selon la plateforme.

\subsection{Formats et finitions possibles}

\begin{center}
\small
\begin{tabularx}{\linewidth}{P{3.4cm}L}
\entete \textbf{Élément} & \textbf{Ce qu’on trouve en général} \\
Formats courants & 11 × 17,5 cm (poche), 12,5 × 19 cm, 13,5 × 21,5 cm
  (roman grand format), 14,8 × 21 cm (A5), 15,24 × 22,86 cm (6 × 9 pouces),
  17 × 24 cm, 21 × 29,7 cm (A4) \\
Reliures & dos carré collé (broché) presque partout ; relié cartonné
  chez certains ; agrafé pour les livrets \\
Papiers intérieurs & blanc ou crème/ivoire, 80 à 90 g/m² ; papier couché
  pour la couleur \\
Couverture & pelliculage brillant ou mat ; rabats, dorure et vernis
  sélectif rarement disponibles \\
Impression & noir (le moins cher) ; couleur standard ou premium, nettement
  plus chère \\
Nombre de pages & souvent de 24 à 700–800 pages ; un nombre pair (parfois
  multiple de 4) est exigé \\
\bottomrule
\end{tabularx}
\end{center}

\subsection{Préparer des fichiers qui passent du premier coup}

\begin{checklist}
  \item \textbf{PDF imprimeur} : PDF/X-1a ou PDF/X-4 de préférence,
        polices incorporées.
  \item \textbf{Fond perdu} de 3 mm (0,125 pouce chez les plateformes
        américaines) si des images ou aplats touchent le bord.
  \item \textbf{Marges intérieures} suffisantes (au moins 15 à 20 mm
        côté reliure pour un livre épais).
  \item \textbf{Images à 300 dpi} à la taille d’impression ; textes noirs
        en noir pur (pas de « noir quadri »).
  \item \textbf{Épaisseur du dos} calculée avec le gabarit de la
        plateforme : elle dépend du nombre de pages \emph{et} du papier.
  \item \textbf{Code-barres EAN-13} de l’ISBN en quatrième de couverture
        (souvent ajouté automatiquement par la plateforme).
  \item \textbf{Page de mentions légales} complète (voir chapitre
        suivant).
\end{checklist}

\begin{astuce}[Gagner du temps]
Les plateformes fournissent des gabarits de couverture gratuits (KDP,
IngramSpark et BoD proposent des générateurs de gabarit en ligne).
Partez toujours de ce gabarit : le dos change à chaque modification du
nombre de pages.
\end{astuce}

\subsection{Distribution : où le livre sera-t-il vendu ?}

C’est la question la plus importante, et celle qui distingue le plus les
plateformes.

\begin{itemize}
  \item \textbf{Boutique de la plateforme} (la plupart en ont une) : utile
        pour la vitrine, mais peu de trafic.
  \item \textbf{Librairies en ligne} : Amazon, Fnac.com, Cultura,
        Decitre, Leslibraires.fr, Place des Libraires… Presque toutes les
        plateformes y donnent accès, avec des délais d’affichage variables.
  \item \textbf{Librairies physiques} : un libraire commande via
        \textbf{Dilicom}, le portail interprofessionnel qui relie libraires
        et distributeurs. Si votre titre n’y est pas référencé avec des
        conditions commerciales correctes (remise, retours), le libraire
        ne pourra pas, ou ne voudra pas, le commander.
\end{itemize}

\begin{attention}[Le piège Amazon]
Un livre publié \emph{uniquement} sur Amazon KDP avec un ISBN gratuit
fourni par Amazon est très difficile à faire commander par une librairie
française : il n’est pas distribué en France hors d’Amazon. Si vous visez
les libraires, prévoyez un second canal (voir chapitre « Mise en place »).
\end{attention}
